\documentclass[12pt]{article}
\usepackage{amsmath, amssymb, amsfonts}
\usepackage[T1]{fontenc}
\usepackage[utf8]{inputenc}
\usepackage{geometry}
\geometry{margin=2.5cm}
\usepackage{lmodern}

\title{Equations of the IS$_d$ Indicator\\
\large Document Singularity Indicator}
\author{Blázquez Ochando, M. · Ovalle Perandones, M.A. · Prieto Gutiérrez, J.J.\\
Universidad Complutense de Madrid}
\date{2025}

\begin{document}
\maketitle

% ────────────────────────────────────────────────────────────────
\section*{Equation (1) — General structure of the indicator}
% ────────────────────────────────────────────────────────────────

\begin{equation}
IS_d \;=\; \alpha \cdot O_d \;+\; (1-\alpha) \cdot I_d
\qquad \alpha \in [0,1]
\label{eq:isd}
\end{equation}

\noindent
Where $O_d$ is the textual originality component, $I_d$ the bibliometric
impact component, and $\alpha$ the weighting parameter.

% ────────────────────────────────────────────────────────────────
\section*{Equation (2) — Normalisation of $O_d$}
% ────────────────────────────────────────────────────────────────

\begin{equation}
O_d \;=\; \min\!\left(\frac{O_d^{raw}}{P_{99}(O^{raw})},\; 1\right)
\label{eq:od}
\end{equation}

\noindent
Where $P_{99}(O^{raw})$ is the 99th percentile of the $O_d^{raw}$ values
in the corpus.

% ────────────────────────────────────────────────────────────────
\section*{Equation (3) — Raw textual originality $O_d^{raw}$}
% ────────────────────────────────────────────────────────────────

\begin{equation}
O_d^{raw} \;=\;
\frac{\displaystyle\sum_{i \,\in\, \Omega_d} S_i \cdot \log\!\left(1 + f_i^{\,doc}\right)}
     {\;\left|\Omega_d\right| \cdot \log\!\!\left(1 + \dfrac{T_d}{\left|\Omega_d\right|}\right) + \varepsilon\;}
\label{eq:odraw}
\end{equation}

\noindent
Where:
\begin{itemize}
  \item $\Omega_d$ = set of unique n-grams present in document $d$
  \item $|\Omega_d|$ = number of unique n-grams (cardinality of the set)
  \item $f_i^{\,doc}$ = absolute frequency of n-gram $i$ in document $d$
  \item $T_d$ = total number of n-grams extracted from document $d$ (with repetition)
  \item $\varepsilon = 10^{-6}$ = numerical smoothing constant
\end{itemize}

% ────────────────────────────────────────────────────────────────
\section*{Equation (4) — Singularity weight $S_i$}
% ────────────────────────────────────────────────────────────────

\begin{equation}
S_i \;=\;
\underbrace{\log\!\left(\frac{C+1}{cd_i + 0{.}5}\right)}_{\text{Robertson IDF}}
\;\cdot\;
\underbrace{\frac{1}{1 + \log\!\left(1 + f_i^{\,corpus}\right)}}_{\text{frequency penalty (BM25)}}
\label{eq:si}
\end{equation}

\noindent
Where:
\begin{itemize}
  \item $C$ = total number of documents in the corpus
  \item $cd_i$ = number of corpus documents containing n-gram $i$
  \item $f_i^{\,corpus}$ = total frequency of n-gram $i$ across the entire corpus
\end{itemize}

% ────────────────────────────────────────────────────────────────
\section*{Equation (5) — Normalised bibliometric impact $I_d$}
% ────────────────────────────────────────────────────────────────

\begin{equation}
I_d \;=\; \frac{\log(1 + C_d)}{\log(1 + C_{max})}
\label{eq:id}
\end{equation}

\noindent
Where $C_d$ is the number of citations received by document $d$ and
$C_{max}$ is the maximum number of citations in the corpus
($C_{max} = 53{,}982$ in calcia.db).
It holds that $I_d \in [0,1]$ exactly: $I_d = 0$ if $C_d = 0$;
$I_d = 1$ if $C_d = C_{max}$.

% ────────────────────────────────────────────────────────────────
\section*{Equation ($\alpha^*$) — Objective calibration criterion}
% ────────────────────────────────────────────────────────────────

\begin{equation}
\alpha^* \;=\; \arg\min_{\alpha \,\in\, [0,1]}
\left|\,\mathrm{skew}\!\left(\alpha \cdot O_d + (1-\alpha)\cdot I_d\right)\right|
\label{eq:alpha}
\end{equation}

\noindent
On calcia.db: $\alpha^* = 0{.}692 \approx 0{.}70$, with $|\mathrm{skew}| = 0{.}00058$
(global minimum). Kendall's $\tau$ of the deciles remains equal to 1.0 for
all $\alpha \in [0;\, 0{.}90]$.

\end{document}